\chapter{Capstone studio}
\chaptertag{Choose a thread, extend it, and explain it}

The capstone is not a miniature research paper. It is a carefully explained mathematical journey: a motivating question, small examples, one substantial proof, and a clear account of how course ideas connect.

\section{Project menu}

\subsection*{A. Rational approximation and Pell equations}
Choose a nonsquare $d$. Compute the continued fraction of $\sqrt d$, identify convergents that solve $x^2-dy^2=\pm1$, and prove at least one recurring pattern. Include a number-line or lattice illustration.

\subsection*{B. Symmetry, orbits, and divisibility}
Generalize necklace counting to length $n$. Classify possible orbit sizes for one composite $n$, explain stabilizers in concrete language, and derive a divisibility identity. M\"obius inversion is an optional extension, not a prerequisite.

\subsection*{C. Pascal's triangle modulo a prime}
Visualize Pascal's triangle modulo $p$. Develop and prove a digit rule (Lucas's theorem is the natural destination), then explain how Kummer's carry rule and the visual pattern encode the same arithmetic.

\subsection*{D. Prime-counting laboratory}
Compare $\pi(x)$ with $x/\log x$, $\operatorname{Li}(x)$, and the Chebyshev bounds over a range your computer can handle. Explain which plots are evidence and which claims are proved. Reconstruct one link in Chebyshev's proof in full detail.

\section{Required deliverables}

\begin{enumerate}
  \item A four-to-six page mathematical exposition.
  \item At least one original diagram or data visualization.
  \item One theorem stated precisely and proved completely.
  \item A one-page ``proof map'' naming the definitions and earlier results used.
  \item A five-minute oral explanation, followed by questions.
\end{enumerate}

\section{Evaluation rubric}

\begin{center}
\small
\begin{tabularx}{\textwidth}{p{.22\textwidth}Xr}
\toprule
Criterion & Evidence of strong work & Weight \\
\midrule
Mathematical correctness & Claims are true, hypotheses are visible, examples are checked, and the main proof closes every logical gap. & 35\% \\
Exposition & The question motivates the definitions; notation is stable; paragraphs explain why steps matter. & 25\% \\
Connection & At least two course strands are linked in a mathematically meaningful way. & 15\% \\
Visual or computational insight & A diagram or computation reveals structure and is interpreted rather than merely displayed. & 15\% \\
Revision and response & The final version addresses peer/instructor feedback and the oral defense shows ownership. & 10\% \\
\bottomrule
\end{tabularx}
\end{center}

\begin{checkpoint}
Before submitting, underline every sentence that makes a mathematical claim. For each underline, ask: Is it a definition, an observed example, a cited theorem, or a claim I have proved? Ambiguous status is often the real source of unclear exposition.
\end{checkpoint}

\section{A final synthesis prompt}

In 500--700 words, respond to this statement:

\begin{quote}
Number theory repeatedly converts multiplication into something easier to organize: remainders, exponent vectors, additive logarithms, or equal-sized orbits.
\end{quote}

Use one precise example from three different chapters. Your response should include at least one equation and explain not only what the conversion does, but why it helps a proof.

